\section{Conclusions and Future Work}
\label{sec:conclusions}

A good precision matrix estimator and a good filtering estimator are not the same thing. On cosmological data, the Bodnar et al.~\cite{bodnar2016direct} methods achieve Frobenius losses half those of Ledoit-Wolf, yet Ledoit-Wolf yields the lowest EnKF RMSE. The reason is stability: a slightly biased estimator that behaves consistently from cycle to cycle accumulates less error over 100 sequential updates than one that is accurate on average but noisy. Without domain-specific prior knowledge, Ledoit-Wolf is the safest default among the methods we tested.

Direct precision shrinkage remains the better choice when the target matches the true precision structure. Our cosmological experiments confirm Looijmans et al.~\cite{looijmans2024comparison}: with near-diagonal data, these methods outperform all alternatives in standalone estimation. Their failure on Lorenz~96 is a consequence of applying diagonal targets to a system whose precision matrix has dense off-diagonal structure, not a deficiency of the framework itself.

The scaled identity target we propose occupies a middle ground---modest gains over the plain identity without requiring domain-specific information, but no substitute for a well-chosen target when one is available. Its ad-hoc factor of 2, borrowed by analogy from the cosmological $2/N_\ell$ term, lacks theoretical justification and could be improved through a data-driven selection criterion.

Several directions remain open. Our implementation does not combine shrinkage with spatial localization~\cite{gaspari1999construction,nino2018ensemble}, which could yield further gains on spatially extended systems. The AR(1) model in the cosmological experiments preserves the covariance structure of the BOSS mock catalogs but is far simpler than a realistic cosmological simulation. Future work should test on higher-dimensional systems such as the quasi-geostrophic model, explore adaptive target selection based on observed covariance structure, and investigate whether the sequential accumulation effect we observed generalizes to other classes of estimators beyond those tested here.

\subsubsection{Acknowledgements.} This work was supported by the Department of Computer Science and Statistics at Universidad del Norte.
